\section*{Abstract}

\textbf{Background:} Stroke remains a leading global cause of death and disability, with an increasing burden driven by aging populations. This study aimed to evaluate one-year stroke outcomes in England by comparing stroke patients to a matched non-stroke cohort. It utilized traditional statistics and explainable machine learning to investigate how patient characteristics, COVID-19 infection and vaccination status, and hyper-acute treatments affect long-term survival and time spent at home.

\textbf{Methods:} We analyzed a nationwide cohort of 388,345 stroke patients (2018–2023) within the NHS England Secure Data Environment, alongside 236,630 age-, sex-, and ethnicity-matched controls. Outcomes evaluated over the first year included survival, time in NHS inpatient care, and home-time (days alive outside inpatient care). The analysis applied descriptive statistics, a Cox proportional hazards model and explainable machine learning algorithms—specifically XGBoost integrated with SHAP (SHapley Additive exPlanations) values—to predict patient outcomes and isolate the relative contribution of specific clinical features.

\textbf{Results:} Compared with matched controls, in the year after stroke, people who had stroke spent 21 more days in inpatient NHS care, 86 fewer days of home-time, and 66 more days dead. The one-year fatality rate for the stroke cohort was 22.0\%. Survival and home-time were most heavily influenced by stroke severity, age, prior disability, and early reperfusion treatment (intravenous thrombolysis and mechanical thrombectomy). Notably, Black and Asian patients demonstrated a survival advantage, exhibiting 65\% and 67\% lower mortality hazard ratios compared to White patients, respectively, despite predicting higher inpatient care time. Unvaccinated COVID-19 infection increased mortality and inpatient time, but prior COVID-19 vaccination largely mitigated the effects of COVID infection.

\textbf{Conclusions:} Stroke imposes a profound and quantifiable deficit on a patient's first year of post-event life. Integrating population-wide registry data with explainable machine learning reveals critical ethnic variations in survival, the protective role of COVID-19 vaccines, and the vital necessity of rapid hyper-acute treatment to optimize patient home-time.